\chapter{Registry Integrity}\label{ch06:registry}

\section{Registries as the kernel's memory}\label{ch06:role}

The kernel components of the preceding chapters are stateless disciplines:
references, boundaries, and encodings constrain values in flight. The claim
and theorem registries are where the program's assertions come to rest. A
\emph{claim registry} is the ledger of empirical hypotheses and product
claims the program tracks, each with its scope, falsifier, evidence
packets, review state, and product-claim status; a \emph{theorem registry}
is the ledger of mathematical statements---axioms, definitions, theorems,
propositions, lemmas, corollaries, conjectures, models---each with its
proof status, migration disposition, novelty status, and admission ring.
Both ledgers are governed by the Project Coherence record contracts, which
fix the field inventory, the enumeration vocabularies, and the conditional
requirements that couple them.

Registries are therefore the highest-value target in the kernel's threat
model. Every downstream consumer that wants to know ``what does the
program claim, and with what status?'' reads a registry; an adversary---or
an honest bug---that can make a registry \emph{appear} to say something it
does not say has subverted the program's memory without touching any other
mechanism. And the historical failure class of Chapter 1 applies with full
force: a ledger row containing the word \textsc{licensed} is precisely the
kind of structural artifact that gets read as authority. The registry
design must accordingly do two jobs at once that pull in opposite
directions: bind the in-memory view to exact bytes so tightly that no
substitution is representable, while ensuring that even a perfectly bound,
perfectly parsed \textsc{licensed} row authorizes nothing.

This chapter develops both jobs. Sections~\ref{ch06:bytebacked} and
\ref{ch06:recheck} give the exact-byte backing and the
reconstruct-and-recheck loading rule, including the constructor-
substitution attack the rule exists to defeat.
Section~\ref{ch06:rejection} walks the full rejection catalog.
Sections~\ref{ch06:immutability} and \ref{ch06:digests} cover recursive
immutability and per-record digests. Section~\ref{ch06:metadata} states
the metadata-is-not-authority invariant, which is where the chapter's two
jobs meet.

\begin{boundary}[Chapter scope]\label{ch06:bd:scope}
As throughout the monograph, statements that the loader ``rejects'' or
``verifies'' describe the v0.1 candidate's tested behavior on the bound
Gate-4 fixture families (registry loader rejection cases are an enumerated
component of the suite). They are bounded software receipts, not
guarantees over untested inputs. The registries themselves, and every view
the kernel constructs of them, are Ring-4 research artifacts: admission
\textsc{open}, product \textsc{not licensed}, disclosure withheld.
\end{boundary}

\section{Exact-byte backing}\label{ch06:bytebacked}

\subsection{Bytes first, views second}

The foundational design decision is about what a loaded registry \emph{is}.
The conventional answer---a parsed data structure, with the source file
forgotten once parsing succeeds---is exactly wrong for evidence, because
it makes the parse the artifact of record. Two different parsers, or two
versions of one parser, can produce different views of the same bytes;
a view can be constructed with no bytes behind it at all; and once the
bytes are gone, no consumer can distinguish a faithful view from a
fabricated one. The kernel inverts this: the artifact of record is the
byte string, and a loaded registry is a \emph{claim about bytes} that
carries the bytes with it.

\begin{definition}[Byte-backed registry]\label{ch06:def:bytebacked}
A loaded registry is a value comprising: the retained raw ledger bytes and
the retained raw schema bytes; the dual artifact identities (raw and
CRLF-to-LF comparison digests with line-ending census,
Chapter 5) computed from each; two externally supplied pinned digests, one
per byte string; the parsed header, schema document, and record tuple; a
canonical digest for the schema document and for each record; and the
kernel research boundary. Construction fails unless every derived
component reconstructs from the retained bytes and every identity matches
its external pin.
\end{definition}

Two clauses of this definition carry the weight. First, the bytes are
\emph{retained}, not merely read: the loaded value holds the exact ledger
and schema bytes for its entire lifetime, so any later verification can go
back to ground truth without touching the filesystem, whose content may
have changed since load. Second, the expected digests are
\emph{externally supplied}: the loader's caller passes the pinned raw
digest of the ledger and of the schema as arguments, and the loader
verifies the bytes it read against those pins before anything else.

\begin{ruleenv}[External pins are trust inputs]\label{ch06:rule:pins}
The pinned digests a registry is loaded against come from the caller---
ultimately from a reviewed manifest outside the loaded artifact---and
never from the loaded bytes themselves. A ledger cannot vouch for its own
identity: a digest stored inside the ledger would be recomputed by
whoever rewrote the ledger, per the self-hash analysis of
Chapter 4. The loader's contract is conditional and honest: \emph{if}
the caller's pin identifies the intended bytes, \emph{then} the loaded
view is bound to those bytes; the kernel takes no position on whether the
pin itself is correct, which is a provenance question outside v0.1's
scope \cite{prov}.
\end{ruleenv}

\subsection{Both ledger and schema are pinned}

It deserves emphasis that the schema document---the machine-readable
record contract the ledger is validated against---is pinned and retained
with exactly the same discipline as the ledger. An unpinned schema is a
laundering channel: an adversary who cannot touch the ledger could
substitute a permissive schema, and previously invalid content would
validate. The v0.1 loader reads the schema bytes, checks them against the
caller's schema pin, validates the schema's own structure (it must be a
closed-object contract: the canonical draft dialect, the canonical
identifier for its kind, additional properties forbidden, and required
and declared fields exactly matching the contract embedded in the kernel),
and only then validates records. The kernel thus checks the schema against
its own embedded expectations as well as against the pin---a deliberate
redundancy, so that neither a wrong pin alone nor a drifted embedded
contract alone can silently widen what validates.

\section{The reconstruct-and-recheck loading rule}\label{ch06:recheck}

\subsection{The rule}

Byte retention makes strong verification \emph{possible}; the loading rule
makes it \emph{mandatory and repeated}.

\begin{ruleenv}[Reconstruct and recheck]\label{ch06:rule:reconstruct}
Every parsed component of a loaded registry---header, schema document, and
each record---is reconstructed from the retained raw bytes and compared,
under canonical encoding, against the component the registry presents.
The full reconstruction runs at construction time and again at every
consuming operation (record resolution, serialization, integrity query).
A registry whose presented components do not reconstruct from its
retained bytes is rejected with a typed raw-binding error naming the
component that failed.
\end{ruleenv}

Concretely, the recheck re-derives the artifact identities from the
retained ledger and schema bytes and compares them field-by-field
(including the line-ending census) against the stored identities and the
external pins; re-decodes the ledger bytes with the full duplicate-key-
rejecting parse; re-validates the root fields, header, and schema
document; re-parses and re-validates every record from the decoded
ledger; recomputes every record's canonical digest; and compares the
freshly reconstructed record tuple against the presented one---kind,
identity, version, canonical payload digest, and boundary, record by
record. Only if every comparison agrees does the operation proceed.

This is the registry-layer instance of Chapter 4's recheck-not-trust
invariant, and the same rationale applies: validation performed once, at
load, protects nothing against in-process mutation afterward, and
protects nothing against values that never went through the load path at
all. The second case is the interesting one, and it is the subject of the
next subsection.

\subsection{The substitution attack}\label{ch06:attack}

The load path---read bytes, verify pins, parse, validate, construct---is
not the only way to obtain a value of the registry type in a memory-safe-
but-not-capability-safe host language. The registry and record types are
ordinary records with public constructors, and the language provides
generic copy-with-changes machinery for such records. This creates the
\emph{substitution attack}: construct a registry value directly, or copy
an honestly loaded one field-by-field, keeping the retained bytes, the
identities, and the pins from the honest artifact while substituting the
parsed content---a record tuple with one extra claim, one altered status
field, one relaxed schema document. The forged value would present
impeccable byte credentials (the bytes \emph{are} the honest bytes) around
a payload the bytes do not contain. Any consumer that trusted the
credentials and read the payload would be reading fiction with a valid
pedigree.

\begin{threat}[Parsed-content substitution]\label{ch06:threat:substitution}
An agent constructs a registry or record value through the public
constructor or generic dataclass-copy machinery, retaining an honest
artifact's raw bytes, identities, and pinned digests while substituting
parsed header, schema, or record content. The presented view diverges
from the retained bytes while every byte-level credential verifies.
\end{threat}

\subsection{The defense}

The defense is that Rule~\ref{ch06:rule:reconstruct} runs \emph{inside
construction itself}, not merely inside the blessed load path. The
registry type's initializer---executed identically whether the value is
built by the loader, by a direct constructor call, or by dataclass-copy
(which routes through the same initializer)---performs the full
reconstruction: it re-parses the retained bytes and compares the result
against whatever parsed content the constructor was handed. The attack is
therefore self-defeating:

\begin{proposition}[Substitution is unrepresentable]
\label{ch06:prop:substitution}
No value of the registry type can simultaneously retain byte strings $B$
and present a parsed component that is not the canonical parse of $B$.
Construction with substituted content fails the raw-binding comparison;
construction with substituted bytes fails the pin comparison; and
post-construction mutation of either side is caught by the same
reconstruction rerun at the next consuming operation. The only values of
the registry type that exist are those whose views are recomputed-and-
matched parses of their own retained, pinned bytes.
\end{proposition}

\begin{designargument}
The initializer receives bytes and parsed components as separate inputs
and derives the comparison itself; there is no constructor parameter that
suppresses the check, and the copy machinery cannot bypass the
initializer. A forger must therefore supply bytes whose canonical parse
equals the forged content---which is not substitution but honest
authorship of a different ledger, and that ledger's raw digest will not
match the external pin the consumer supplies. The residual path,
in-memory mutation after construction, does not create a stable forged
value: every consuming operation reruns the reconstruction
(Rule~\ref{ch06:rule:reconstruct}), so the mutated value fails at first
use. As always, an adversary who rewrites the kernel's own code is out of
scope (Chapter 4).
\end{designargument}

The corollary for reviewers is a pleasant inversion of the usual audit
burden: to check what a v0.1 registry value says, it suffices to check
what its retained bytes say, because no representable value of the type
can disagree with its bytes. The bytes are reviewable with independent
tools; the binding is the kernel's job and is exercised by the bound
fixtures; the meaning of the bytes is the record contract's job, to which
we now turn.

\subsection{The load procedure}

For reference in the rejection catalog, the complete load path is worth
setting out as an ordered procedure; each step names the class of
rejection it can raise, and the catalog of Section~\ref{ch06:rejection}
expands each class.

\begin{protocol}[Registry loading]\label{ch06:prot:load}
Loading a registry of a declared kind, against caller-supplied expected
raw digests for ledger and schema, proceeds:
\begin{enumerate}
\item resolve the kind to its embedded contract (canonical schema-version
string, record-schema name, record-identifier grammar, field
inventories); an unsupported kind fails closed;
\item validate both caller-supplied pins as well-formed digest values
before any input is read, so a malformed pin cannot degrade into a
skipped comparison;
\item read the exact ledger bytes; compute the dual artifact identity;
reject on unreadable input, non-UTF-8 content, duplicate JSON keys,
or a non-object root;
\item compare the ledger's raw digest against its pin; mismatch fails
closed before any content-level validation runs, so no diagnostic
about unpinned content is ever emitted;
\item validate the root field set exactly; validate the embedded schema
identifiers against the kind's contract; extract and validate the
header;
\item read the schema bytes (from the declared companion location or a
caller-specified path); compute their identity; compare against the
schema pin; validate the schema document as the closed-object
contract for the kind;
\item parse and validate every record against the field inventories,
shape rules, enumeration vocabularies, identifier grammar, and
conditional contracts; compute each record's canonical digest;
construct each record value (whose own initializer revalidates);
\item construct the registry value, whose initializer enforces kind
homogeneity, identity uniqueness, boundary exactness, and then runs
the full reconstruction of Rule~\ref{ch06:rule:reconstruct} against
the retained bytes before the value ever escapes to the caller.
\end{enumerate}
\end{protocol}

The step order embodies a policy: identity before content. The pins are
checked before any semantic validation, so every error message the
loader can produce is a statement about bytes the caller intended to
load; the loader never leaks a detailed parse diagnosis of an artifact
that was not the pinned one, which would otherwise be a convenient
oracle for probing unpinned content through error channels.

\section{The rejection catalog}\label{ch06:rejection}

The loader's validation is deliberately exhaustive, and this section
enumerates it. The organizing principle is a closed-world reading of the
record contracts: everything not explicitly admitted is rejected, every
rejection carries a typed error code naming the violated rule, and no
rejection is downgraded to a warning. The catalog is presented in the
order the loader encounters each class.

\subsection{Byte and parse layer}

\emph{Non-retained or non-UTF-8 bytes.} The decode layer requires exact
byte strings (not paths, not streams, not text already decoded by an
unknown convention) and rejects content that is not valid UTF-8. This
anchors everything downstream: there is exactly one byte string in play,
and its decoding is unambiguous.

\emph{Duplicate JSON keys.} The parse installs a pair-collision hook so
that an object with a repeated key is rejected at parse time. Standard
JSON parsers silently keep the last binding, which is a smuggling
channel: a ledger containing
\texttt{"product\_claim\_status"} twice---once benign, once elevated---
would validate against schema and read differently depending on parser
convention. The kernel refuses to let parser convention decide evidence
content; a duplicate key is an error, full stop.

\emph{Non-object roots.} A ledger whose top level is not a JSON object is
rejected before any field logic runs.

\subsection{Root and header layer}

\emph{Unknown or missing root fields.} The ledger root must contain
exactly the six canonical fields---schema version, record-schema name,
document status, as-of date, completeness statement, and the record
array---no more and no fewer. The error for a mismatch reports both the
missing and the surplus field names. Surplus rejection is the closed-world
principle at the root: an unrecognized root field is not ``extra
information'' but a place where content invisible to the contract could
ride along inside pinned bytes.

\emph{Wrong schema identifiers.} The root's schema-version and
record-schema fields must equal, exactly, the canonical strings for the
registry's kind (claim ledger versus theorem ledger, each with its own
versioned identifier). A ledger of one kind cannot load as the other, and
a ledger authored against a different contract version fails loudly
rather than best-effort parsing.

\emph{Malformed header.} The header triple extracted from the root is
validated on its own: document status and completeness must be nonblank
strings of admissible canonical text, and the as-of field must be a real
calendar date in ISO form---pattern-checked \emph{and} calendar-checked,
so \texttt{2026-02-30} fails even though it matches the pattern.

\subsection{Schema-document layer}

As Section~\ref{ch06:bytebacked} described, the schema document must
declare the canonical JSON-Schema dialect, carry the exact canonical
identifier for its kind, describe a closed object (additional properties
forbidden), declare required and optional fields exactly matching the
kernel's embedded contract (duplicate entries in the required list are
rejected), and contain its conditional-contract clauses. A schema
document that is valid JSON Schema but not \emph{this} contract is
rejected.

\subsection{Record layer: shape}

\emph{Unknown and missing record fields.} Each record is checked against
the exact required and optional field inventories for its kind---sixteen
required fields for a claim (identity, version, title, statement kind,
claim text, source references, scope, controls, falsifier, uncertainty,
evidence packets, review state, product-claim status, admission ring,
disclosure status, reviewers), seventeen for a theorem---with the same
two-sided error reporting as the root: missing fields are named, surplus
fields are named, and either is fatal.

\emph{Wrong shapes.} Every field has an exact shape requirement: nonblank
strings where prose is required; arrays of strings for reference lists
(with per-kind minimums---at least one source reference always); the
claim scope must be an object with exactly the five canonical subfields
(population, models, architectures, contexts, time window) of the exact
required types; timestamp fields must be RFC 3339 date-times \emph{with
an explicit timezone}, pattern-checked and parse-checked. Shape errors
name the record and field.

\emph{Wrong enumerations.} Every status field is validated against its
closed vocabulary: claim statement kinds (empirical hypothesis or product
claim); review states; product-claim status (\textsc{licensed},
\textsc{not licensed}, or \textsc{open}); disclosure status (open, held,
or flagged for human review); theorem statement kinds; the ten migration
dispositions; the nine proof statuses; the five novelty statuses. A value
outside the vocabulary---including a case variant or a padded string---is
rejected. The admission ring is checked as an exact integer in the
admitted range for its kind (claims: 3 or 4; theorems: 1 through 4), with
the host language's boolean-integer conflation explicitly excluded, so
\texttt{true} cannot pose as ring 1.

\emph{Conditional contracts.} The contracts couple fields across a
record, and the loader enforces the couplings: a claim whose product
status reads \textsc{licensed} must also be in review state
\textsc{reviewed}, at admission ring 3, with nonempty evidence packets,
nonempty reviewers, and a present review timestamp---a \textsc{licensed}
row missing any of these is not a weaker license but a contract
violation, rejected outright; a theorem marked formally verified must be
reviewed and cite proof artifacts; a theorem admitted at ring 1 or 2 must
be reviewed and dated. These conditions encode, at the data layer, the
same necessary-condition discipline as Chapter 4's boundary rules: the
loader ensures that elevated statuses cannot be \emph{written} without
their structural prerequisites, while conferring nothing when they are
(Section~\ref{ch06:metadata}).

\subsection{Closed vocabularies and migration audit}

The theorem registry's vocabularies deserve a paragraph beyond their
role as rejection triggers, because their structure encodes the
program's migration posture. The \emph{migration disposition} vocabulary
distinguishes ten outcomes for a legacy statement under review---
unassessed, in review, inherited, clarified, corrected, conditional,
downgraded, deprecated, refuted, replaced---and the point of the closed
enumeration is that a migrating statement must land in exactly one of
them, explicitly. The historically dangerous alternatives are the open
ones: free-text dispositions, where ``mostly carried over'' and
``carried over'' blur; or absence, where a statement's survival into a
new version is inferred from its continued presence in a document. The
same logic governs the \emph{proof status} vocabulary, which separates
gradations that overclaim tends to collapse---an informal proof present
is not proved; a standard result cited is not a new result proved; a
counterexample present coexists in the vocabulary with proof statuses
precisely so that refutation is a recordable state rather than a
deletion; and \emph{independent replication required} exists so that
``proved here'' and ``established'' remain distinguishable. The
\emph{novelty status} vocabulary does the same for the claim of
originality, separating standard results and standard specializations
from framework couplings and candidate new results, with
novelty-unassessed as the honest default. None of these vocabularies is
enforced semantically by the kernel---the loader cannot check that a
proof exists, only that the status field is a member of the closed set
and that the conditional contracts hold (a formally verified entry must
cite artifacts and be reviewed). The kernel's contribution is
structural: every migration question the program has historically
answered vaguely now has a field that must be filled from a finite
list, and a review of the registry is a review of those explicit
answers.

\subsection{Record layer: identity}

\emph{Invalid identifiers.} Record identifiers must match the declared
identifier grammar for their kind---a fixed uppercase prefix
(\texttt{CLM-} for claims, \texttt{THM-} for theorems) followed by an
alphanumeric-initial body over an explicit uppercase alphabet. The
grammar is a machine-identifier profile in the sense of Chapter 5:
ASCII, closed alphabet, no invisible characters, no confusables.

\emph{Identity--payload agreement.} The record value's identity and
version fields must equal the corresponding fields inside its payload;
a record cannot present one identity while carrying another.

\emph{Duplicate and ambiguous identities.} Within a registry, the
(identifier, version) pair must be unique; a duplicate is rejected at
load. Resolution is correspondingly strict: resolving an identifier that
matches no record is a not-found error, and resolving an identifier that
matches multiple versions without specifying which is an
\emph{ambiguity} error, not a default to latest. ``Latest'' is a
convention, and conventions are where readers and writers diverge
silently; the kernel requires the caller to say which version of a claim
it means whenever the question has more than one answer.

\emph{Missing metadata.} Finally, the absence-of-metadata cases: required
fields absent (caught by the inventory check), required timestamps absent
where a conditional contract demands them, empty reference arrays where
minimums apply, and blank strings where nonblank prose is required. The
loader treats missing metadata identically to wrong metadata---rejection
with a named field---because an evidence ledger with holes is not a
smaller ledger but an unvalidatable one.

\begin{remark}[Why warnings are absent]\label{ch06:rem:nowarnings}
The catalog contains no warning tier: every violation is fatal to the
load. The rationale is the desensitization argument of Chapter 5 applied
to validation: a warning tier trains consumers to proceed past
violations, and the one violation that matters will arrive dressed like
the hundred that did not. A registry either satisfies its contract
exactly or is not a registry; partial credit is a property of the
authoring workflow, which can run the loader and read the typed errors,
not of the loaded artifact.
\end{remark}

\subsection{The catalog at a glance}

Table~\ref{ch06:tab:catalog} summarizes the catalog by rejection family
and the layer at which each fires. Every family carries a typed error
class whose code names the family (the implementation prefixes these
\texttt{E-REGISTRY-}, for example
\texttt{E-REGISTRY-DUPLICATE-KEY} for repeated JSON keys and
\texttt{E-REGISTRY-AMBIGUOUS} for underspecified version resolution),
so that fixtures and operational logs can assert the family, not merely
the fact, of a failure.

\begin{table}[ht]
\centering
\begin{tabular}{lll}
\toprule
Layer & Rejection family & Representative trigger \\
\midrule
Bytes & unreadable / non-bytes input & path unreadable; non-byte payload \\
Bytes & encoding & ledger not valid UTF-8 \\
Parse & duplicate JSON key & repeated key in any object \\
Parse & root shape & top level not an object \\
Pin & ledger hash mismatch & raw digest differs from pin \\
Pin & schema hash mismatch & schema digest differs from pin \\
Root & root field set & missing or surplus root field \\
Root & schema identifiers & wrong kind or contract version \\
Header & header shape / date & blank status; invalid as-of date \\
Schema & contract mismatch & open object; field set drift \\
Record & field inventory & missing or unknown record field \\
Record & shape & wrong type; malformed scope; bad timestamp \\
Record & enumeration & value outside closed vocabulary \\
Record & ring & non-integer or out-of-range ring \\
Record & conditional contract & licensed row without reviewers \\
Record & identifier grammar & id fails prefix/alphabet rule \\
Record & identity agreement & id/version differ from payload \\
Registry & duplicate identity & repeated (id, version) pair \\
Registry & kind homogeneity & foreign-kind record in registry \\
Registry & boundary & non-Ring-4 kernel boundary \\
Resolve & not found & no record for cited id \\
Resolve & ambiguity & multiple versions, none specified \\
Binding & raw reconstruction & view diverges from retained bytes \\
\bottomrule
\end{tabular}
\caption{The registry rejection catalog by layer. Each family is fatal;
there is no warning tier. ``Binding'' failures are raised by the
reconstruction rule at construction and at every consuming operation.}
\label{ch06:tab:catalog}
\end{table}

\subsection{Worked rejections}

Three of the families repay concrete illustration, because the inputs
that trigger them look innocent.

\begin{example}[Duplicate-key smuggling]\label{ch06:ex:dupkey}
A crafted claim record contains, in its serialized bytes, the pair
\texttt{"product\_claim\_status":"NOT LICENSED"} followed later in the
same object by \texttt{"product\_claim\_status":"LICENSED"}. Under a
last-binding parser the record reads as licensed; under a first-binding
parser, as not licensed; a validator using one convention and a
consumer using the other disagree about the same pinned bytes---the
worst possible property for evidence. The kernel's parse hook rejects
the object at the first repeated key, before validation begins, so the
bytes are simply not a loadable ledger under any convention.
\end{example}

\begin{example}[Ambiguity is not latest]\label{ch06:ex:ambiguous}
A ledger legitimately carries two versions of one claim---say versions
\texttt{1.0} and \texttt{1.1} of \texttt{CLM-EXAMPLE}---as the record
contracts permit. A consumer resolves \texttt{CLM-EXAMPLE} without a
version, expecting ``the current one.'' The kernel raises the ambiguity
error instead of choosing, because the two plausible conventions
(latest by version-string order, latest by review date) can disagree
with each other and with the author's intent, and a citation minted
under one convention would verify differently under the other. The
consumer must cite \texttt{CLM-EXAMPLE} at an exact version; the
citation quadruple of Definition~\ref{ch06:def:citation} then names one
record forever.
\end{example}

\begin{example}[Schema swap]\label{ch06:ex:schemaswap}
An agent leaves a pinned ledger untouched but supplies a widened schema
document---additional properties permitted, an enumeration extended---
hoping previously rejected content will validate in a future load of a
modified ledger. The swap fails twice before it can matter: the schema
bytes fail their own external pin (step 6 of
Protocol~\ref{ch06:prot:load}), and even a schema that somehow carried
a correct pin is validated against the kernel's \emph{embedded}
contract, which rejects an open object or a drifted field set outright.
The double check means a schema-layer attack must subvert the pin
manifest and the kernel code together, not either alone.
\end{example}

\section{Recursive immutability}\label{ch06:immutability}

A loaded registry is shared: multiple components hold references to it,
and its records are handed out on resolution. If the parsed values were
ordinary mutable structures, every recipient could alter what every other
recipient sees---an aliasing channel that would undermine the byte
binding from inside, since Rule~\ref{ch06:rule:reconstruct} catches
mutation only at the next consuming operation, and an in-place mutation
between operations could mislead a reader that consulted the structure
directly.

The kernel therefore freezes loaded values recursively at construction.
Every parsed mapping is wrapped in a read-only mapping view; every parsed
array is converted to an immutable sequence carrier that preserves
JSON-array semantics; and the freeze descends through the entire
structure, so there is no depth at which a mutable container survives.
Scalars are validated as they are frozen---strings must be admissible
canonical text (the control- and format-character rules of Chapter 5
apply to ledger content as to everything else), and the reserved
canonical type tag is rejected wherever it appears in ledger data, which
closes the tag-forgery channel of Chapter 5 at the registry layer as well
as at the encoder.

Two design details deserve note. First, the immutable array carrier is a
distinct type, not the host language's tuple: the canonical encoding
distinguishes tuples from JSON arrays (Chapter 5), so an array frozen
into a tuple would change its canonical digest---immutability must not
alter identity. The freeze machinery accordingly rejects bare tuples in
registry data as a shape error rather than guessing which identity was
meant. Second, thawing---recovering plain structures for revalidation---
is the exact inverse and is used only inside the reconstruction rule,
where the recovered structure is immediately compared against a fresh
parse of the retained bytes; frozen and thawed forms are two views of one
canonical content, with equality of canonical digests as the invariant
connecting them.

\begin{invariant}[Frozen views, stable digests]\label{ch06:inv:frozen}
The recursive freeze applied at load, and the thaw applied during
reconstruction, preserve canonical content: the canonical digest of a
record's payload is identical before freezing, after freezing, and after
a freeze--thaw round trip. Immutability changes who can write; it never
changes what the bytes say.
\end{invariant}

\section{Per-record digests bound to registry bytes}\label{ch06:digests}

Consumers rarely cite a whole registry; they cite a record---\emph{this}
claim, at \emph{this} version. The kernel therefore gives each record its
own deterministic digest: the canonical-encoding digest (Chapter 5) of
the record's complete payload, computed from content that
Rule~\ref{ch06:rule:reconstruct} has bound to the retained ledger bytes.
The record digest inherits the byte binding transitively: it is not an
independent assertion about the record but a compressed name for content
whose provenance is the pinned ledger.

\begin{definition}[Record citation]\label{ch06:def:citation}
A record citation is the quadruple (registry kind, record identifier,
record version, record canonical digest), optionally accompanied by the
registry's pinned raw digest. A consumer holding a citation can verify,
against any independently obtained copy of the ledger bytes, that the
cited record exists in those bytes with exactly the cited content---
without trusting the kernel instance that issued the citation.
\end{definition}

This is the registry layer's contribution to the program's evidence
chain: downstream artifacts (compatibility receipts, evidence packets,
the runtime gate's decision records) refer to claims by citation, and
the citation is checkable end-to-end---ledger pin to ledger bytes, bytes
to parse, parse to record, record to digest---with each link exercised
by the reconstruction rule. Per the standing distinction, the whole
chain establishes tamper evidence and reference integrity only. A
correct citation of a \textsc{licensed} record proves that the ledger
says \textsc{licensed}; what that is worth is the next section's
question.

\subsection{Serialized views re-verify}

The registry and its records export serialized views---for display, for
inclusion in receipts, for cross-process transfer---and the export path
is disciplined the same way as every other consuming operation: a
record's view export first re-runs the record's full integrity check,
and a registry's view export first re-runs the complete reconstruction
of Rule~\ref{ch06:rule:reconstruct}. An exported view therefore cannot
be produced from a value that has drifted from its bytes; the last act
before serialization is a fresh proof of binding. The exported view
itself is ordinary canonical-encodable content (Chapter 5), so a
consumer can digest it, and the record's embedded canonical digest
gives the consumer a cross-check that survives the export boundary:
recompute the digest of the exported payload and compare it to the
digest the export carries. The pattern generalizes: every time registry
content crosses a boundary---load, resolve, export---the crossing is
preceded by reverification and followed by content that carries the
means of its own recheck. What the pattern never produces, per the next
section, is content that carries the means of its own authorization.

\section{Two registries, one discipline}\label{ch06:tworegistries}

The division into a claim registry and a theorem registry is itself a
design decision worth a section, because it enforces at the ledger
layer a distinction the program has historically blurred: the
difference between an empirical assertion about systems in the world
and a mathematical assertion about defined structures. The two kinds
have different falsification semantics (a claim names its falsifier
and its uncertainty as required fields; a theorem names its
assumptions, dependencies, and proof artifacts), different status
vocabularies (product-claim status and disclosure status on one side;
proof status, migration disposition, and novelty status on the other),
and different admission-ring ranges (claims are admitted only at rings
3 and 4---an empirical claim is never ``reviewed mathematics''---while
theorems span rings 1 through 4). A single mixed ledger would invite
category slippage: an empirical result drifting into citation as
established theory, or a theorem's proof status being read as evidence
for an empirical deployment claim. Separate kinds, separate contracts,
separate identifier prefixes, and a loader that rejects a record of
one kind in a registry of the other make the slippage a shape error.

The kinds are nonetheless linked, deliberately and in one direction:
a claim record may carry theorem references among its optional fields,
citing the mathematical statements it depends on, and a theorem record
carries its dependencies on other theorems. Both are citation lists
over the identifier grammars---strings, validated for shape, not
resolved by the loader. The v0.1 kernel does not chase references
across registries, does not verify that a cited theorem exists, and
does not compute transitive dependency closures; those are integrity
properties of the \emph{corpus}, not of a single pinned artifact, and
checking them requires fixing which set of registries constitutes the
corpus---exactly the trusted-manifest question that
Rule~\ref{ch06:rule:pins} defers to the process layer. The kernel's
per-artifact discipline is designed so that a future corpus-level
checker has sound material to work with: every reference it would
chase is grammatically valid, every target it would resolve is
byte-pinned and citable by digest, and every version it would need is
explicit because ambiguity is unresolvable
(Example~\ref{ch06:ex:ambiguous}). Cross-registry closure is recorded
as an open obligation in Chapter 12.

\section{Metadata is not authority}\label{ch06:metadata}

\subsection{The invariant}

Everything in this chapter so far makes registries harder to falsify.
Nothing in it makes registries \emph{authoritative}, and the design goes
out of its way to keep the two properties separate.

\begin{invariant}[Metadata is not authority]\label{ch06:inv:metadata}
The text \textsc{licensed} inside a registry record is necessary
metadata, never sufficient authority. Loading, validating, digesting,
and citing a record whose product-claim status reads \textsc{licensed}
establishes only that the pinned ledger bytes carry a contractually
well-formed row with that status. No kernel operation treats that fact
as authorization for any claim-bearing action; in v0.1 the runtime
gate's answer to claim-bearing intent is constant denial regardless of
registry content, because the kernel has no authenticated reviewer and
no legal trust root against which a \textsc{licensed} row could be
verified as an act rather than a string.
\end{invariant}

The invariant's phrasing---necessary, never sufficient---is exact. The
conditional contracts of Section~\ref{ch06:rejection} make well-formed
\textsc{licensed} metadata a \emph{necessary} condition for any future
authorization: a future reviewed design will require, at minimum, a
reviewed, ring-3, evidence-backed, reviewer-named, dated record, and the
loader guarantees that nothing less can even be written. But the
sufficiency direction requires what v0.1 deliberately lacks: a trust
root. The reviewers field of a record is a list of names---strings,
authored by whoever authored the ledger. Nothing binds a name to a
person, a person to a review act, or a review act to an entitlement.
Verifying such bindings is the function of authentication
infrastructure---signatures, enrolled identities, revocation---which is
a separate technology from integrity checking and cannot be improvised
out of it \cite{anderson}. A kernel that pretended otherwise would be
manufacturing the historical failure this monograph opened with:
structurally impeccable receipts read as authority.

\subsection{Descriptive queries are marked as such}

The kernel does provide a convenience query that reports whether a claim
record's payload satisfies the full licensed-metadata contract
(status \textsc{licensed}, state \textsc{reviewed}, ring 3, nonempty
packets and reviewers, dated). Its existence is worth a remark because
of what it is careful \emph{not} to be: it is a descriptive predicate
over ledger content, documented as such, returning a fact about bytes;
it feeds display and reporting, and no decision path consumes it as a
permission. The runtime gate (Chapter 9) does not call it to decide
anything, because the gate's decision on claim-bearing intent is not a
function of ledger content at all in v0.1---it is the constant
\emph{no}. Keeping the descriptive query and the decision path visibly
disconnected is a small design choice with a large purpose: the day a
future design connects them, the connection must be an explicit,
reviewed edge binding a trust root, not the quiet reuse of a boolean
that was lying around.

\subsection{Worked scenario: the perfect forged ledger}

The chapter's disciplines compose into a useful worst-case exercise.
Suppose an agent authors, from scratch, a complete claim ledger in which
every record is contractually perfect and every status reads
\textsc{licensed}: correct fields, valid enumerations, ring 3,
\textsc{reviewed}, plausible reviewer names, dated, with evidence-packet
identifiers that resolve to artifacts the agent also authored. Every
mechanism in this chapter behaves as designed: the bytes parse; the
records validate; digests bind; citations verify. What has the agent
obtained?

Within the kernel: nothing. The loaded registry is a Ring-4 open
artifact like every other (its kernel boundary is pinned to Ring-4 open
regardless of what its records say---the layering of Chapter 4); the
descriptive query truthfully reports that the ledger says
\textsc{licensed}; and the runtime gate denies claim-bearing intent
exactly as it would against an empty ledger, because no code path
consults ledger content for authorization. Outside the kernel: the
agent still needs a caller to load the forged ledger, which requires
the caller to supply its digest as a pin---and pins, by
Rule~\ref{ch06:rule:pins}, come from reviewed manifests, not from the
artifact's author. The forgery thus fails twice over: it cannot pass
the pin gate without subverting the manifest process, and even loaded,
it authorizes nothing. The residual risk is the one the kernel cannot
reach and does not claim to: a human reader shown the forged ledger
outside any kernel discipline may believe it. The kernel's contribution
is that belief would have to be formed \emph{against} the kernel's
every output, not through it.

\section{Limits and obligations}\label{ch06:limits}

The chapter closes by naming what the registry design leaves open; the
items join the obligation set of Chapter 12.

First, the pin-provenance gap: the loader verifies bytes against
caller-supplied pins, and the discipline governing the pins themselves
---the reviewed manifest, its own integrity, its update process---is a
process-layer obligation outside the v0.1 kernel. The registry
machinery is exactly as trustworthy as the weakest pin its caller
supplies.

Second, the authority gap, stated throughout: no authenticated reviewer
identity, no legal trust root, no authorization semantics for
\textsc{licensed} rows. The future-design constraints (an exact trusted
registry and schema manifest among them) are recorded with the runtime
decision model in Chapter 9.

Third, the coverage bound: every ``rejects'' in this chapter is
evidenced by the bound fixture families of the Gate-4 receipt, which
exercise representatives of each rejection class, not the space of all
malformed inputs. The reconstruct-and-recheck rule reduces what a
consumer must trust to the correctness of one reconstruction function;
it does not prove that function correct, and independent review of it
remains an open obligation.

\begin{checkclass}[Registry fixture families]\label{ch06:check:fixtures}
The bound suite's registry families track
Table~\ref{ch06:tab:catalog}: (i) pin-mismatch fixtures for ledger and
schema; (ii) parse-layer fixtures (duplicate keys, non-UTF-8 bytes,
non-object roots); (iii) root, header, and schema-contract fixtures
covering missing, surplus, and drifted fields; (iv) per-field record
fixtures for each shape rule, each closed vocabulary, the identifier
grammar, and each conditional contract, for both claim and theorem
kinds; (v) identity fixtures (payload disagreement, duplicate and
ambiguous identities, not-found and underspecified resolution); (vi)
substitution fixtures constructing registry and record values with
swapped parsed content, swapped bytes, and subclassed types, asserting
the typed binding failures of
Proposition~\ref{ch06:prop:substitution}; and (vii) immutability and
digest fixtures asserting frozen views, freeze--thaw digest stability
(Invariant~\ref{ch06:inv:frozen}), and citation verification against
independently re-read bytes. As throughout, passing receipts for these
families are bounded historical evidence, not conformance claims.
\end{checkclass}

Within those limits, the registry layer completes the kernel's
structural story: bytes are pinned, views are reconstructed and
rechecked, every irregularity in the catalog is a typed rejection,
loaded values are frozen, records are citable by digest---and at the
end of all that machinery, the word \textsc{licensed} is still just a
word. The next chapters take up the execution side of the kernel: the
hash-verified legacy compiler profile (Chapter 7), its concurrency and
failure disciplines (Chapter 8), and the decision model (Chapter 9)
where the constant-false authorization invariant turns this chapter's
distinction between metadata and authority into an enforced runtime
answer.
